\documentclass[../../../include/open-logic-section]{subfiles}

\begin{document}

\olfileid{sth}{spine}{Valphabasic}
\olsection{层级的基本性质}

为了揭示 $V_\alpha$ 定义的奠基性意义，我们将介绍关于它们的一些基本结果。首先给出一个定义：\footnote{“有能的”并非标准术语；这是 \citet{ButtonLT1} 使用的名称。}

\begin{defn}
	集合 $A$ 是\emph{有能的}，当且仅当 $\forall x((\exists y \in A)x \subseteq y \lif x \in A)$。 
\end{defn}
\begin{lem}\ollabel{Valphabasicprops}
对每个序数 $\alpha$：
\begin{enumerate}
	\item\ollabel{Valphatrans} 每个 $V_\alpha$ 都是传递的。
	\item\ollabel{Valphapotent} 每个 $V_\alpha$ 都是有能的。
	\item\ollabel{Valphacum} 若 $\gamma \in \alpha$，则 $V_\gamma
	\in V_\alpha$（从而由 \olref{Valphatrans} 亦有 $V_\gamma \subseteq V_\alpha$）。
\end{enumerate}
\end{lem}

\begin{proof}
我们通过（同时）超限归纳法来证明。归纳假设：对每个序数 $\beta < \alpha$，\olref{Valphatrans}--\olref{Valphacum} 成立。 

当 $\alpha = \emptyset$ 时是平凡的。 

设 $\alpha = \ordsucc{\beta}$。要证 \olref{Valphacum}：若
$\gamma \in \alpha$，则由归纳假设有 $V_\gamma \subseteq V_\beta$，
故 $V_\gamma \in \Pow{V_\beta} = V_\alpha$。要证
\olref{Valphapotent}：设 $A \subseteq B \in V_\alpha$，即 $A
\subseteq B \subseteq V_\beta$，则 $A \subseteq V_\beta$，于是 $A \in
V_\alpha$。要证 \olref{Valphatrans}：若 $x \in A \in
V_\alpha$，我们有 $A \subseteq V_\beta$，故 $x \in V_\beta$；而由归纳假设 $V_\beta$ 是传递的，所以 $x
\subseteq V_\beta$，从而 $x \in V_\alpha$。 

设 $\alpha$ 是极限序数。要证 \olref{Valphacum}：若
$\gamma \in \alpha$，则 $\gamma \in \ordsucc{\gamma} \in \alpha$，于是由归纳假设 $V_\gamma \in V_{\ordsucc{\gamma}}$，从而
$V_\gamma \in \bigcup_{\beta \in \alpha} V_\beta = V_\alpha$。要证
\olref{Valphatrans} 和 \olref{Valphapotent}：只需注意到传递集合的并集仍是传递的，有能集合的并集仍是有能的。 
\end{proof}

\begin{lem}\ollabel{Valphanotref}
对每个序数 $\alpha$，$V_\alpha \notin V_\alpha$。
\end{lem}

\begin{proof}
由超限归纳法。显然 $V_\emptyset \notin V_\emptyset$。 

若 $V_{\ordsucc{\alpha}} \in V_{\ordsucc{\alpha}} = \Pow{V_\alpha}$，
则 $V_{\ordsucc{\alpha}} \subseteq V_\alpha$；而由 \olref{Valphabasicprops} 知 $V_\alpha
\in V_{\ordsucc{\alpha}}$，即得 $V_\alpha \in V_\alpha$。反之，若 $V_\alpha \notin V_\alpha$
则 $V_{\ordsucc{\alpha}} \notin V_{\ordsucc{\alpha}}$。

若 $\alpha$ 是极限序数且 $V_\alpha \in V_\alpha = \bigcup_{\beta \in
\alpha}V_\beta$，则存在 $\beta \in \alpha$ 使得 $V_\alpha \in V_\beta$；但此时亦有 $V_\beta \in V_\alpha$，故由 \olref{Valphabasicprops}（两次）得 $V_\beta \in
V_\beta$。反之，若对所有 $\beta \in \alpha$ 均有 $V_\beta \notin
V_\beta$，则 $V_\alpha \notin V_\alpha$。
\end{proof}

\begin{cor}
对任意序数 $\alpha, \beta$：$\alpha \in \beta$ 当且仅当 $V_\alpha \in V_\beta$。
\end{cor}

\begin{proof}
\Olref{Valphabasicprops} 给出了其中一个方向。反之，假设 $V_\alpha \in V_\beta$。则由 \olref{Valphanotref} 知 $\alpha \neq \beta$；且 $\beta \notin \alpha$，否则将有 $V_\beta \in V_\alpha$，从而由 \olref{Valphabasicprops}（两次）得 $V_\beta \in V_\beta$，与 \olref{Valphanotref} 矛盾。故由三分律得 $\alpha \in \beta$。
\end{proof}
\noindent
以上结果使我们得以将每个 $V_\alpha$ 视为该层级的第 $\alpha$th
阶段。原因如下。

显然，我们的 $V_\alpha$ 可以看作是在一个\emph{迭代}过程中形成的，因为我们对序数的使用恰好反映了迭代的概念。此外，若一个阶段在另一个阶段之前形成，即 $V_\beta \in V_\alpha$，亦即 $\beta \in \alpha$，则我们的形成过程是\emph{累积的}，因为 $V_\beta \subseteq V_\alpha$。最后，我们确实构成了在先前任何阶段中可用的\emph{所有}集合的汇集，因为任何后继阶段
$V_{\ordsucc{\alpha}}$ 都是其前驱 $V_\alpha$ 的幂集。

简而言之：借助 $\ZFminus$，我们\emph{几乎}已经阐明了集合的累积迭代层级这一图景。（当然，我们仍需证成替换公理模式。）

\end{document}